%% ma: L10-B14-0067
%% lop: 10
%% chuong: 5
%% bai: 14
%% dang: 10.14.8
%% loai: TN4
%% muc: VD
%% dap_an: "A"
%% nguon: "chuyen-de-cac-so-dac-trung-cua-mau-so-lieu-khong-ghep-nhom-toan-10-chuong-trinh-moi.pdf (GV. Phan Nhật Linh) (P05), Bài 03 Dạng 2 Phần I Câu 11"
%% kien_thuc: "Lớp 10 Bài 13 (số trung bình), Bài 14 (phương sai, độ lệch chuẩn)"
%% trang_thai: chinh-thuc
%% ghi_chu: "Tự duyệt theo yêu cầu giáo viên 22/10/2026."
\begin{ex}
Để khảo sát kết quả môn Toán (thang điểm 10) trong kì thi tuyển sinh năm vừa qua của trường A, người ta chọn một mẫu gồm $100$ học sinh. Điểm môn Toán của các học sinh được cho ở bảng tần số sau:
\begin{center}
\begin{tabular}{|c|c|c|c|c|c|c|c|c|c|c|c|}\hline
Điểm & $0$ & $1$ & $2$ & $3$ & $4$ & $5$ & $6$ & $7$ & $8$ & $9$ & $10$\\\hline
Tần số & $1$ & $1$ & $3$ & $5$ & $8$ & $13$ & $19$ & $24$ & $14$ & $10$ & $2$\\\hline
\end{tabular}
\end{center}
Phương sai của mẫu số liệu (làm tròn đến hàng phần trăm) là
\choice{\True $3{,}96$}{$3{,}99$}{$3{,}98$}{$3{,}97$}
\loigiai{$\sum n_ix_i=623$; $\sum n_ix_i^2=4277$.

$s^2=\dfrac{4277}{100}-\left(\dfrac{623}{100}\right)^2=3{,}9571\approx3{,}96$. Chọn A.}
\end{ex}
